\section{Introduction}
\label{sec:introduction}
% P:introduction-01
In an aqueous electrolyte that splits into two liquids or exchanges solvent with a vapor, phase splitting, ion partitioning and chemical speciation determine the observables that a parameter fit compares with data and that a process model constrains.
A mean ionic activity coefficient, a phase composition, a bubble pressure or the pH of a liquid is then a property of a solved equilibrium, not of a state at given composition.
Least-squares estimation needs the derivatives of such observables with respect to the parameters, and an equation-oriented process model needs their derivatives with respect to temperature, pressure and feed, including the response of every solved phase and species.
The accuracy of these derivatives governs the Jacobian-based diagnostics and the convergence of equation-oriented models~\citep{Allan2024}.

% P:introduction-02
The perturbed-chain statistical associating fluid theory (PC-SAFT) of Gross and Sadowski~\citep{Gross2001,Gross2002} describes the residual Helmholtz energy of chain molecules with hard-chain, dispersion and association contributions.
Electrolyte PC-SAFT (ePC-SAFT) adds a Debye--H\"uckel term for ion--ion interactions~\citep{Cameretti2005,Held2008,Held2014}, and later versions add a Born term for ion solvation and a composition-dependent relative permittivity~\citep{Bulow2019,Bulow2020,Bulow2021,Figiel2025}.
These equations are taken here from their sources; none of them is new.

% P:introduction-03
Several open-source packages implement PC-SAFT with exact Helmholtz-energy derivatives.
teqp evaluates the Helmholtz energy of an equation of state and obtains its derivatives by automatic differentiation~\citep{bellImplementingEquationState2022}.
FeOs derives the properties of its models from one Helmholtz-energy implementation through generalized dual numbers and provides parameter estimation and classical density functional theory~\citep{rehnerFeOsOpenSourceFramework2023}.
Clapeyron.jl organizes a large family of equations of state around a shared set of property and equilibrium methods with automatic differentiation~\citep{walkerClapeyronjlExtensibleOpenSource2022}.
CoolProp evaluates multiparameter reference equations of state for pure and pseudo-pure fluids and tabulates check values for independent implementers~\citep{bellPurePseudopureFluid2014}; where such a reference equation exists it is the more accurate source of pure-fluid properties, and this package does not replace it.
ThermoPack provides cubic and PC-SAFT equations of state in one thermodynamic framework, whose recommended citation reviews the established methods such frameworks rest on~\citep{wilhelmsenThermodynamicModelingEquations2017}.
At the versions surveyed, Clapeyron.jl differentiates solved nonreactive equilibria implicitly, FeOs differentiates bubble pressures with respect to parameters, and Clapeyron.jl solves electroneutral liquid--liquid equilibria with ions in both liquids (Table~\ref{tab:software}).

% P:introduction-04
The equilibrium methods are also established.
Constrained Gibbs-energy minimization over phases and species is a general formulation of phase and chemical equilibrium~\citep{Pereira2012,Tsanas2019,Coatleven2022}, and interior-point solution of such problems has been reported~\citep{Belov2023}.
Simultaneous ePC-SAFT reaction and phase equilibrium has been published for acid dissociation in liquid--liquid systems~\citep{Ascani2022pH} and for esterification~\citep{Ascani2023}, and Tsanas et al.\ solved reactive phase equilibria with ions in several liquids using an electrolyte nonrandom two-liquid model~\citep{Tsanas2022}.
Tangent-plane analysis tests the stability of a candidate phase set~\citep{Michelsen1982}, and differentiation of the optimality conditions of a nonlinear program gives the sensitivity of its solution~\citep{pirnayOptimalSensitivityBased2012}.

% P:introduction-05
This work presents the ePC-SAFT package, distinguished here from the published ePC-SAFT equations it evaluates.
Its contribution is the demonstrated combination of these elements in one native calculation (Fig.~\ref{fig:workflow}).
A parameter record selects one electrolyte Helmholtz energy, and a single Gibbs-energy-minimization driver solves neutral, charged and reacting phase-equilibrium specifications for it on one set of residual equations.
Differentiation of those equations gives exact derivatives of the solved state, which serve as Jacobian columns of parameter estimation and as first and second derivatives of an equation-oriented process model.
No equivalent combination was identified in the sources surveyed at the versions named in Table~\ref{tab:software}; this is a survey finding with unknown entries, not a statement about every implementation.
The combination holds within the tested phase supports, the sets of species allowed in each phase, and within the reference conventions and derivative orders stated in Sections~\ref{sec:equilibrium} and~\ref{sec:derivatives}.
No overall speed advantage is claimed; Section~\ref{sec:performance} defines the operation-specific computational-cost comparison.
The measured-data comparisons of Section~\ref{sec:physical} state the resulting physical accuracy and its failures.

\begin{figure}[pos=H]
\centering
% F1: request flow through the one native calculation (schematic, no data).
\begin{tikzpicture}[
  font=\footnotesize,
  box/.style={cse node, text width=30mm, minimum height=13mm},
  wide/.style={cse node, text width=37mm, minimum height=19mm, fill=cseSky!10},
  search/.style={cse node, text width=34mm, minimum height=11mm, dashed},
  flow/.style={cse arrow},
  node distance=6mm and 7mm]
  \node[box] (record) {Parameter record\\contributions, $\epsilon_r$ relation};
  \node[box, below=of record] (thermal) {Thermal records,\\reactions, $K(T)$};
  \node[wide, right=of $(record.east)!0.5!(thermal.east)$] (energy) {Helmholtz energy, Eq.~(\ref{eq:residual})\\recorded once, derivatives of any order;\\association and density closures};
  \node[box, right=12mm of energy, yshift=17mm] (state) {State properties\\$P$, $\mu_i$, $h$, $c_p$ and actions};
  \node[box, right=12mm of energy] (equil) {Declared equilibrium\\$\mathbf{F}(\mathbf{y},\boldsymbol\theta)=\mathbf{0}$\\Newton attempt, then Ipopt};
  \node[box, below=of equil] (actions) {Solved-state actions\\$\mathbf{F}_\mathbf{y}\,\mathbf{y}_\theta=-\mathbf{F}_\theta$};
  \node[search, below=10mm of energy, xshift=-8mm] (search) {Explicit phase search:\\tangent plane, discovery;\\on explicit request only};
  \node[box, right=of state, yshift=-6mm] (fit) {Parameter estimation\\least squares, minimax};
  \node[box, right=of actions, yshift=6mm] (process) {Process models\\Pyomo/IDAES via ASL};
  \draw[flow] (record.east) -- (energy.west |- record.east);
  \draw[flow] (thermal.east) -- (energy.west |- thermal.east);
  \draw[flow] (energy.east) -- ++(4.5mm,0) |- (state.west);
  \draw[flow] (energy.east) -- (equil.west);
  \draw[flow] (equil.south) -- (actions.north);
  \draw[flow] (state.east) -- (fit.west |- state.east);
  \draw[flow] (actions.east) -- (process.west |- actions.east);
  \draw[flow] (actions.north east) -- (fit.south west);
  \draw[flow, dashed] (search.east) -- (search.east -| {$(energy.east)+(8.5mm,0)$}) |- ([yshift=-4mm]equil.west);
  \path ([yshift=-3mm]current bounding box.south) node[inner sep=0] {};
\end{tikzpicture}

\caption{One native calculation serves every request.
A parameter record and declared thermal and reaction data define the Helmholtz energy of Eq.~\eqref{eq:residual} (Section~\ref{sec:model}), evaluated with its algorithmic derivatives on the selected support and thermal interval.
State properties, the rows of a declared equilibrium (Section~\ref{sec:equilibrium}) and the directional derivatives, called actions, of its solved state (Section~\ref{sec:derivatives}) are derivatives of that energy; parameter fits (Section~\ref{sec:estimation}) and process models consume them.
An explicit phase search uses the same driver but is never run inside a declared solve.
ASL, AMPL Solver Library.}
\label{fig:workflow}
\end{figure}



% P:introduction-06
Each comparison in this work answers one of four questions, and every figure panel is labelled with its class.
Numerical verification asks whether the code solves the stated equations; it compares results with independent codes, analytic identities and step refinement (Section~\ref{sec:verification}).
Published-model reproduction re-evaluates a published formulation and its parameters.
Comparison with calculated source values tests numerical translation; comparison with measured data or evaluated correlations instead assesses that published model while retaining its original fitting overlap and source limitations.
Calibration reports how well fitted parameters represent the data used to fit them, and prediction reports the error of fixed parameters on data outside the fit (Sections~\ref{sec:estimation} and~\ref{sec:physical}).
Inherited literature parameters may have been fitted to some of the data shown, so a prediction with such parameters is labelled with that overlap.
No comparison here is called physical validation, which would additionally require independent measurements with uncertainties, judged against a defined use.


% P:introduction-07
The examples span pure-fluid saturation, neutral and associating mixtures, electrolytes, solvent transfer and reacting phase splits.
The propane fit connects calibration to withheld-temperature prediction; mixed salts expose failed transfer from fitted binaries.
Sections~\ref{sec:model}--\ref{sec:estimation} introduce equations and algorithms, Section~\ref{sec:implementation} their use, and Sections~\ref{sec:physical} and~\ref{sec:verification} physical and numerical accuracy.
Cost, engineering use and limitations follow in Sections~\ref{sec:performance}--\ref{sec:conclusions}.


\begin{table}[pos=htbp]
\centering
\caption{Equilibrium and derivative capabilities of five PC-SAFT implementations, from their source code and documentation.
G, one general formulation shared with other specifications; S, a dedicated routine; P, partial or restricted scope; U, assembled by the user from lower-level outputs; X, excluded by the source or documentation; ?, not shown in the sources checked, which is unknown rather than absent.
The codes describe how a capability is provided, not its accuracy or cost.
An asterisk marks a class in which a comparison in Section~\ref{sec:verification} computes the same case with that code and this package.
Versions: FeOs~0.9.5 source, with 0.10.1 for the derivative, pressure--enthalpy and pressure--entropy flash and fitting classes; Clapeyron.jl development source following release~0.6.20, with the 0.6.29 documentation for implicit derivatives and the generalized flash; teqp~0.23.2; ThermoPack development source.
Supporting Information Section~\ref{sec:si-sources} gives the source of each cell.
\textsuperscript{a}Ions confined to the aqueous liquid or at trace amounts in the organic liquid.
\textsuperscript{b}Internal-energy--volume specifications and pressure--enthalpy or pressure--entropy specifications with charged species are not available.
\textsuperscript{c}Phase envelopes are verified for neutral vapor--liquid systems.
\textsuperscript{d}Derivative orders depend on the case and the input direction (Section~\ref{sec:derivatives}).
\textsuperscript{e}No physical comparison is reported in this work.}
\label{tab:software}
\small
\begin{tabularx}{\linewidth}{@{}Xccccc@{}}
\toprule
Problem class & This work & Clapeyron.jl & FeOs & teqp & ThermoPack \\
\midrule
Phase envelopes and critical points & G\textsuperscript{c} & S & S* & S* & S \\
Tangent-plane stability analysis & G & S & S* & ? & S \\
Neutral liquid--liquid and three-phase equilibrium & G & S* & S* & S & P \\
Electrolyte activity and osmotic coefficients & S & S* & U* & ? & X \\
Born solvation and composition-dependent permittivity & S & S* & S* & ? & X \\
Electrolyte vapor--liquid equilibrium with speciation & G\textsuperscript{e} & X & ? & ? & X \\
Electrolyte liquid--liquid equilibrium, ions in both liquids & G & S* & ? & ? & X \\
Simultaneous reaction and phase equilibrium & G & X & ? & ? & ? \\
Reactive electrolyte liquid--liquid equilibrium & G\textsuperscript{a} & X & ? & ? & X \\
Caloric properties and pressure--enthalpy flash & G\textsuperscript{b} & G & S* & G & S \\
Exact derivatives of equilibrium results & G\textsuperscript{d} & P* & P* & P & P \\
Electrolyte and reactive parameter estimation & S & P & ? & ? & ? \\
\bottomrule
\end{tabularx}
\end{table}
